\documentclass[aspectratio=169, 14pt]{beamer}
\usepackage{beamer_preamble}

\begin{document}

% ==== Title ===============================================================
\titleslide{Lesson 3-5 \textperiodcentered\ Period 3}{Real + Complex Zeros, and Modeling}

% ==== Do Now ==============================================================
\begin{frame}{Do Now --- Intersection Points}{Algebra 2 \textperiodcentered\ Lesson 3-5 \textperiodcentered\ Period 3}
\phasebadges{1--2}{5}{bridge from Tuesday}
\vspace{0.2cm}
\begin{projcard}{Practice (3-5-savvas-q11) \textperiodcentered\ Where do the graphs cross?}{slidegold}
At what points do the graphs of
\(f(x) = x^3 - 2x^2 - 16x + 20\) and \(g(x) = -12\) intersect?

\vspace{0.3cm}
\small\textcolor{slidegray}{Graph is on your packet. List all intersection points.}

\vspace{0.2cm}
\textbf{Bridge (verbal):} A cubic has zeros \(2\), \(1+\sqrt{3}\), and \(1-\sqrt{3}\).
Are the coefficients rational, irrational, or complex? Why?
\end{projcard}
\end{frame}

% ==== Launch ==============================================================
\begin{frame}{Launch --- Complex Zeros + Modeling Rules}{Algebra 2 \textperiodcentered\ Lesson 3-5 \textperiodcentered\ Period 3}
\phasebadges{2}{12}{teacher models}
\vspace{0.2cm}
\begin{projcard}{Complex-Zero Rules (keep on your page)}{slidegreen}
\begin{enumerate}
  \setlength{\itemsep}{0pt}
  \item Complex zeros of a real-coefficient poly come in \textbf{CONJUGATE PAIRS}.
  \item If \(a+bi\) is a zero, then \(a-bi\) must also be a zero.
  \item A degree-\(n\) polynomial has exactly \(n\) zeros (counting multiplicity).
  \item Example: cubic with real zero \(x = 4\) and complex zero \(2+3i \Rightarrow\) third zero is \(2-3i\).
\end{enumerate}
\end{projcard}
\vspace{-0.1cm}
\begin{projcard}{Modeling Rules (Storage Box)}{slidegreen}
\begin{enumerate}
  \setlength{\itemsep}{0pt}
  \item Write \(V(x)\) in factored form; each factor = a physical dimension.
  \item Use the constraint (``length is greatest'') to label length, width, height.
  \item To find dimensions at a given volume, set up and solve the cubic.
\end{enumerate}
\end{projcard}
\end{frame}

% ==== Explore =============================================================
\begin{frame}{Explore --- TPS + DOK-3 Modeling}{Algebra 2 \textperiodcentered\ Lesson 3-5 \textperiodcentered\ Period 3}
\phasebadges{2--3}{38}{student work}
\small\textcolor{slidegray}{4 items in order. Plan \(\sim\)20 min for Practice \#27 (Storage Box).}

\vspace{0.15cm}
\begin{projcard}{Try It 3a \textperiodcentered\ Real + complex zeros (cubic)}{slideblue}
\(f(x) = 2x^3 - 8x^2 + 9x - 9\). Find ALL zeros.
\end{projcard}
\vspace{-0.15cm}
\begin{projcard}{Try It 3b \textperiodcentered\ Real + complex zeros (quartic)}{slideblue}
\(f(x) = x^4 - 3x^2 - 4\). Find ALL zeros.
\end{projcard}
\vspace{-0.15cm}
\begin{projcard}{Practice \#17 \textperiodcentered\ Real + complex zeros}{slideblue}
\(f(x) = x^3 - x^2 - 4x - 6\). Find ALL zeros.
\end{projcard}
\vspace{-0.15cm}
\begin{projcard}{Practice \#27 \textperiodcentered\ STORAGE BOX MODELING (\(\bigstar\))}{slideaccent}
\(f(x) = x^3 + 2x^2 - 3x\) = volume; \(x\) = width (ft). Parts (a)--(d): factor, zeros, label dimensions, find dims when \(V = 10\ \text{ft}^3\).
\end{projcard}
\end{frame}

% ==== Share / Summary =====================================================
\begin{frame}{Share / Summary}{Algebra 2 \textperiodcentered\ Lesson 3-5 \textperiodcentered\ Period 3}
\phasebadges{2}{7}{DOK-3 reveal + assessment preview}
\vspace{0.2cm}
\begin{projcard}{Sentence frame \textperiodcentered\ Storage Box}{slideblue}
``For the Storage Box, the factored form was \underline{\hspace{1.6cm}}.
I chose \underline{\hspace{1cm}} as the length because \underline{\hspace{1.8cm}}.
To find \(V = 10\ \text{ft}^3\), I \underline{\hspace{2.2cm}}.''
\end{projcard}

\vspace{0.15cm}
\begin{projcard}{Connection to Thursday's Assessment}{slideaccent}
Assessment Q\#6 (Lucy's tray) uses the \textbf{same move}: factored volume model \(\rightarrow\) extract a hidden dimension.\\
Today's Storage Box is your rehearsal.
\end{projcard}
\end{frame}

% ==== Exit Ticket =========================================================
\begin{frame}{Exit Ticket --- Summary Recap}{Algebra 2 \textperiodcentered\ Lesson 3-5 \textperiodcentered\ Period 3}
\phasebadges{1--2}{3}{not graded}
\vspace{0.2cm}
\begin{projcard}{Complete on your exit ticket}{slidegold}
\begin{enumerate}
  \setlength{\itemsep}{0.2cm}
  \item Today I learned that \underline{\hspace{5cm}}.
  \item For the Storage Box, the hardest step was \underline{\hspace{4.5cm}}.
  \item Thursday's assessment will test \underline{\hspace{4.5cm}}.
\end{enumerate}
\end{projcard}
\end{frame}

\end{document}
